\documentclass[11pt,a4paper]{article}
% Common preamble for the hint sheets (\input from each file)
\usepackage[utf8]{inputenc}
\usepackage[T1]{fontenc}
\usepackage{lmodern}
\usepackage{amsmath,amssymb,amsthm,mathtools}
\usepackage[a4paper,margin=2.3cm]{geometry}
\usepackage{enumitem}
\usepackage{xcolor}
\usepackage[hidelinks]{hyperref}
\renewcommand{\leq}{\leqslant}
\renewcommand{\le}{\leqslant}
\renewcommand{\geq}{\geqslant}
\renewcommand{\ge}{\geqslant}
\setlength{\parindent}{0pt}
\setlength{\parskip}{0.35em}
% \hintstitle{Sheet title}{Class work / Homework}
\newcommand{\hintstitle}[2]{\begin{center}{\LARGE\bfseries Hints: #1\par}\vspace{0.5em}{\large #2\par}\end{center}
\noindent\textit{The hints to each problem are cumulative. Read one, try again, and open the next only if you are still stuck.}\par\medskip\hrule\medskip}
\newcommand{\problem}[1]{\par\bigskip{\large\bfseries Problem #1}\par\nopagebreak\smallskip}
\newcommand{\subpart}[1]{\par\smallskip\textbf{(#1)}\par\nopagebreak}
\newcommand{\hint}[1]{\par\smallskip\noindent\textcolor{black!70}{\textbf{Hint #1.}}\enspace}
\newcommand{\note}[1]{\par\smallskip\noindent\textit{Note.} #1\par}
\providecommand{\lt}{<}
\providecommand{\gt}{>}
\providecommand{\tightlist}{\setlength{\itemsep}{0pt}\setlength{\parskip}{0pt}}
\emergencystretch=3em

\begin{document}
\hintstitle{Euclidean geometry -- 2. Conics, areas and projective maps}{Homework}

\problem{1}
\hint{1} The affine map $x=3u$, $y=v$ turns the ellipse into the unit circle and multiplies all areas by the same factor $3$. So the two areas are equal iff the images are equal.
\hint{2} The images are circular sectors. The line $y=\frac12x$ becomes $v=\frac32u$, which makes the angle $\theta_1=\arctan\frac32$ with the $u$-axis; the line $y=mx$ becomes $v=3mu$, and the second sector has the angle $\frac\pi2-\arctan(3m)$.
\hint{3} Equal sectors: $\arctan(3m)=\frac\pi2-\arctan\frac32$, i.e.\ $3m=\frac23$. Answer: $m=\frac29$.

\problem{2}
\hint{1} The standard parabola with vertex $(p,q)$ is $y=(x-p)^2+q$. Write down the condition ``the point $(x_0,y_0)$ lies on the standard parabola with vertex $(p,q)$'' and look at it from the other side.
\hint{2} $y_0=(x_0-p)^2+q$ is equivalent to $-q=(p-x_0)^2-y_0$, i.e.\ to ``$s(V)=(p,-q)$ lies on the standard parabola with vertex $s(A)=(x_0,-y_0)$''.
\hint{3} $A_i$ lies on the parabolas with vertices $V_j$ and $V_k$, so $s(V_j)$ and $s(V_k)$ lie on the parabola with vertex $s(A_i)$. Hence the parabolas with vertices $s(A_j)$ and $s(A_k)$ both pass through $s(V_i)$; two different standard parabolas meet in at most one point (their difference is linear), so $s(V_i)$ is their intersection point.

\problem{3}
\hint{1} An ellipse can only come from a plane $z=\alpha x+\beta y+\gamma$ that is not too steep: $\alpha^2+\beta^2<1$, $\gamma>0$. On the section $z=\sqrt{x^2+y^2}$.
\hint{2} So the projection is the curve $\sqrt{x^2+y^2}=\alpha x+\beta y+\gamma$ in the plane $Oxy$. The right-hand side is $\sqrt{\alpha^2+\beta^2}$ times the distance from $(x,y)$ to the line $l\colon\alpha x+\beta y+\gamma=0$ (the curve lies on the side where this expression is positive).
\hint{3} Thus $|XO|=e\,d(X,l)$ with $e=\sqrt{\alpha^2+\beta^2}<1$: an ellipse with focus $O$ and directrix $l$.

\problem{4}
\hint{1} Lemma: if the tangent at a point $X$ of the ellipse meets the directrix $l$ corresponding to $F$ at a point $Y$, then $\angle XFY=90^\circ$.
\hint{2} Check the lemma in coordinates: $F=(c,0)$, $l\colon x=\frac{a^2}c$, tangent at $X=(x_0,y_0)$: $\frac{xx_0}{a^2}+\frac{yy_0}{b^2}=1$. Then $Y=\bigl(\frac{a^2}c,\frac{b^2}{y_0}(1-\frac{x_0}c)\bigr)$ and $\overrightarrow{FX}\cdot\overrightarrow{FY}=(x_0-c)\frac{b^2}c+b^2\bigl(1-\frac{x_0}c\bigr)=0$.
\hint{3} Let the tangents at $A$ and $B$ meet $l$ at $Y_A$ and $Y_B$. Both $FY_A$ and $FY_B$ are perpendicular to the line $AFB$ at $F$, so they are the same line, and $Y_A=Y_B$ is its intersection with $l$. Hence $C=Y_A=Y_B$ lies on the directrix and $CF\perp AB$.

\problem{5}
\hint{1} Put the ellipse in canonical position, $O=(u,v)$. On the ellipse $y^2=b^2-\frac{b^2}{a^2}x^2$, so $|XO|^2=e^2x^2-2ux-2vy+b^2+u^2+v^2$. The circle touches the ellipse exactly where this function attains its maximum $r^2$ on the ellipse.
\hint{2} The conic $r^2-\bigl(e^2x^2-2ux-2vy+b^2+u^2+v^2\bigr)=0$ has no $xy$ and no $y^2$ terms and touches the ellipse at the two points of contact, so it lies in the pencil spanned by the ellipse and the double line through these points. Comparing the coefficients of $xy$ and $y^2$ shows that $u=0$ or $v=0$: $O$ lies on an axis.
\hint{3} If $v=0$, the function $e^2x^2-2ux$ is convex in $x\in[-a,a]$, so the maximum is at an end of the major axis, a single point unless $u=0$; then $O$ is the centre, $r=a$, $OF_1=c$. If $u=0$, maximise $e^2a^2+b^2+v^2-\frac{c^2}{b^2}y^2-2vy$ over $y$: the value is $r^2=\frac{a^2}{c^2}(c^2+v^2)$. In both cases $e^2r^2=c^2+v^2=OF_1^2$.

\problem{6}
\hint{1} The maps $(x,y)\mapsto(tx,\frac yt)$ preserve areas and the hyperbola. Alternatively, compute: for $A=(a,\frac1a)$ the tangent at $P=(p,\frac1p)$ has slope $-\frac1{p^2}$ and $AB$ has slope $-\frac1{ab}$.
\hint{2} The area of $APB$ is maximal when the tangent at $P$ is parallel to $AB$ (Theorem 3), i.e.\ $p=\sqrt{ab}$.
\hint{3} The map with $t=\frac pa=\frac bp$ sends $A$ to $P$ and $P$ to $B$, hence the first region onto the second. (Or: the region between the chord from $(a,\frac1a)$ to $(p,\frac1p)$ and the hyperbola has area $\frac12\bigl(\frac pa-\frac ap\bigr)-\ln\frac pa$, a function of $\frac pa$ only.)

\problem{7}
\hint{1} By an affine map the parabola becomes $y=x^2$; area ratios do not change. Let $A,B$ have abscissas $\alpha<\beta$ and the $P_i$ abscissas $\alpha=x_0<x_1<\dots<x_n<x_{n+1}=\beta$.
\hint{2} The polygon is the parabolic segment over $AB$ minus the $n+1$ small segments over the sides, so its area is $\frac{(\beta-\alpha)^3}6-\frac16\sum_{i=0}^n(x_{i+1}-x_i)^3$ (Theorem 3).
\hint{3} With $d_i=x_{i+1}-x_i$ and $\sum d_i=\beta-\alpha$, the sum $\sum d_i^3$ is minimal exactly for equal $d_i$ (power means or convexity), and then equals $\frac{(\beta-\alpha)^3}{(n+1)^2}$. So the maximum exists, it is attained at equally spaced abscissas, and the ratio is $1-\frac1{(n+1)^2}=\frac{n(n+2)}{(n+1)^2}$.

\problem{8}
\hint{1} Take the parabola $y=x^2$ (focus $(0,\frac14)$, directrix $y=-\frac14$) and use the abscissas $a,b,c$ as parameters. The tangent at $t$ is $y=2tx-t^2$, and the tangents at $s$ and $t$ meet at $\bigl(\frac{s+t}2,st\bigr)$.
\hint{2} (a) The foot of the perpendicular from the focus to the tangent at $t$ is $(\frac t2,0)$: all three feet lie on the tangent at the vertex, so by Simson's theorem the focus lies on the circumcircle of $ABC$. (b) The altitude from $C$ is perpendicular to the tangent at $c$; it meets $y=-\frac14$ at the point with abscissa $\frac{a+b+c}2+2abc$, which is symmetric in $a,b,c$.
\hint{3} (c) $S_{abc}=\frac12|(b-a)(c-a)(c-b)|$, and the same determinant for $A,B,C$ gives $S_{ABC}=\frac14|(a-b)(b-c)(c-a)|$. (d) $S_{abC}=\frac14|b-a|^3$, and similarly for the others; since $b$ lies between $a$ and $c$, $|c-a|=|c-b|+|b-a|$.

\problem{9}
\hint{1} A projective map of the plane sends lines to lines and preserves the cross-ratio of four lines through a point (it is the cross-ratio of their intersections with any line, Theorem 4). It sends a non-degenerate conic to a circle, so it suffices to prove the statement for a circle.
\hint{2} On a circle the angles $\angle AXB$, $\angle BXC,\dots$ do not depend on $X$ (inscribed angles, oriented modulo $\pi$), and the cross-ratio of four lines through $X$ is expressed through the sines of these angles.
\hint{3} Pascal for $ABCDEF$: let $P=AB\cap DE$, $Q=BC\cap EF$, $R=CD\cap FA$. Project $B,D,E,F$ from $A$ onto the line $DE$ and from $C$ onto the line $EF$. You get the points $P,D,E,M$ ($M=AF\cap DE$) and $Q,N,E,F$ ($N=CD\cap EF$) with equal cross-ratios.
\hint{4} So there is a projective map from the line $DE$ to the line $EF$ sending $P,D,E,M$ to $Q,N,E,F$. It fixes the common point $E$, hence it is a perspectivity, and its centre lies on all the lines $PQ$, $DN=CD$, $MF=AF$. Thus $PQ$ passes through $R=CD\cap AF$.

\problem{10}
\hint{1} Each turn is by $120^\circ$, so the segments have only three directions, and segment $i$ is parallel to segment $i+3$. Put the broken line inside an equilateral triangle $T$ with sides parallel to the three directions, and let $x_i$ be the distance from segment $A_iA_{i+1}$ to the side of $T$ parallel to it (indices modulo $3n$).
\hint{2} Parallel segments do not cross. Every crossing is between some segment $i$ and some segment $j+1$ with $j\equiv i\pmod3$, $j\ne i$, i.e.\ it is a crossing of the two-segment paths $A_iA_{i+1}A_{i+2}$ and $A_jA_{j+1}A_{j+2}$. Since the turning is always the same, along every segment the distance to the ``next'' side changes in the same direction; work out that such a crossing is possible only if
\[(*)\qquad (x_i-x_j)(x_{i+1}-x_{j+1})<0,\]
and that the two paths cross at most once. There are $3\binom n2$ unordered pairs $\{i,j\}$ with $i\equiv j\pmod3$.
\hint{3} It remains to find $\bigl\lfloor\frac{n-1}2\bigr\rfloor$ pairs violating $(*)$, since $3\binom n2-\bigl\lfloor\frac{n-1}2\bigr\rfloor=\bigl\lfloor\frac32n^2-2n+1\bigr\rfloor$. Let $x_0$ be the smallest of all $x_i$ and fix $k$ with $1\le k<\frac n2$. Show that one of the pairs $(-6k+j,\,j)$, $j=0,1,\dots,6k-1$, violates $(*)$.
\hint{4} If all of them satisfy $(*)$: $x_{-6k}>x_0$ (no three vertices are collinear, so two parallel segments are at different distances), then $(*)$ gives $x_{-6k+1}<x_1$, then $x_{-6k+2}>x_2$, and so on, alternating; after $6k$ steps $x_0>x_{6k}$, contradicting the minimality of $x_0$. Different $k$ give different pairs (for $k+k'=\frac n2$ the two families are disjoint because $j$ stops at $6k-1$), so this gives $\bigl\lfloor\frac{n-1}2\bigr\rfloor$ violations.
\end{document}
